\chapter[Die Schuld des Unschuldigen]{Die Schuld des Unschuldigen und die Schuldlosigkeit der Schuldigen}
\quellseite{1}

„Er wechselt mit uns wunderlich“ — sangen wir vor wenigen Wochen von dem Geheimniß der Krippe Christi. Das erfüllte uns mit freudigem Staunen, daß der Herr ein Knecht ward, damit die Knechte Herren würden. „Er wechselt mit uns wunderlich“ — dürfen wir jetzt abermal singen von dem Geheimnisse des Kreuzes Christi. Und wenn auch ein Ton schmerzlicher Klage unsern Gesang diesmal durchzittert, so ist unser freudiges Staunen über den „wunderlichen Wechsel“ gewiß ebenso groß unter dem Kreuz wie an der Krippe. Ja, ist das nicht noch wunderlicher, daß der Heilige ein Sünder wird, damit die Sünder Gerechte werden? „Wie wunderbahrlich ist doch diese Strafe!“

Christus ist der Heilige, nicht bloß der heilige Gott, in dem kein Unrecht ist, dem von Ewigkeit die Seraphim singen, wie sie dem Vater und dem Heiligen Geiste singen: „Heilig, heilig, heilig ist Gott“ — sondern Christus ist auch der heilige Mensch. Der Mensch in seiner Unschuld, wie er rechtschaffen aus der Hand des Schöpfers hervorgegangen war, wandelt wieder auf Erden in Christo. Es war „ein Heiliges“, das von der Maria geboren ward. Und dieses Heilige, dieser heilige Mensch, hat seine Sündlosigkeit bewahrt. Wiewohl allenthalben versucht wie andere Menschen, ist er doch ohne Sünde geblieben. Ein „frommer Mensch“ starb am Fluchholz; ein „Heiliger“ lag im Grabe Josephs, ohne die Verwesung zu sehen.

Was die Schrift von der Unschuld Christi sagt, das sollen wir auf seine Menschheit beziehen. Es wäre doch seltsam, wenn uns die heiligen Männer Gottes wieder und wieder sagen müßten, Gott sei heilig. Der Gott der Christen ist ja nicht den Götzen der Heiden gleich, von denen sich die Heiden allerlei Schande und Greuel zu erzählen pflegten. Es bedarf also in Bezug auf den Gott der Christen nicht erst der Versicherung, daß er heilig sei. Nein, mit der bloßen Erklärung: „Christus ist der wahrhaftige Gott“, ist auch bereits das bezeugt, daß er vollkommen heilig ist.

Aber das ist das Wunderbare, das Unerhörte, daß aus dem ganz verderbten Menschengeschlecht heraus ein untadeliger, unschuldiger, heiliger Mensch erstanden ist.

„Er hat keine Sünde gethan“, 1 Petr. 2, 22.; „er wußte von keiner Sünde“, 2 Cor. 5, 21.; „es ist keine Sünde in ihm“, 1 Joh. 3, 5.; „er hat niemand Unrecht gethan“, Jes. 53, 9.; „kein Betrug ist in seinem Munde erfunden“, 1 Petr. 2, 22. — so und mit vielen andern Worten, in immer neuen Redewendungen bezeugt die Schrift die völlige Unschuld dessen, der selbst seine Zeitgenossen herausforderte, ihn, wenn sie könnten, „einer Sünde zu zeihen“, Joh. 8, 46. Die Schrift will eben jeden Gedanken an die Sündhaftigkeit Christi abschneiden, nicht den geringsten Zweifel an seiner Unschuld aufkommen lassen, darum wiederholt sie ihre Aussage über diese Sache so oft und wechselt doch den Ausdruck. Wollte jemand zwar zugeben, Christus sei in seinem Thun untadelig gewesen, aber nicht also in seinen Worten, so erklärt ihm die Schrift, daß auch in seinem Munde kein Falsch gewesen sei. Und gibt der Zweifler das zu, behält sich aber vor zu glauben, daß Christus mindestens sündhafte Regungen in seinem Herzen verspürt habe, so tritt die Schrift ihm abermals mit einem entschiedenen Nein entgegen. Bis in die geheimsten Gedanken, bis auf den tiefsten Herzensgrund erstreckt sich Christi Unschuld.

Ist es nicht auch wunderbar, wie sich gerade in den letzten Tagen seines Erdenlebens die Zeugnisse für Christi Unschuld, und zwar aus dem Munde und aus den Handlungen seiner Feinde, häufen? Judas bekennt verzweifelnd: „Ich habe übel gethan, daß ich unschuldig Blut verrathen habe“, Matth. 27, 4. Hätten die Hohenpriester eine gerechte Sache wider Christum gehabt, so hätten sie gewiß keine falschen Zeugen gedungen. Selbst mit dem Hülfsmittel der Lüge konnten sie nichts auf Christum bringen, Matth. 26, 60. 61. Pilatus erklärte den Juden einmal über das andere, er finde „keine Schuld“, Joh. 18, 38. 19, 4. 6., „keine Ursache“, Luc. 23, 4., „der Sachen keine“, der er beschuldigt war, V. 14., an Jesu. Herodes auch nicht, denn es war
\quellseite{2}
vor ihm „nichts auf Jesum gebracht, das des Todes werth war“, Luc. 23, 15. Pilatus’ Weib warnt ihren Gatten, er möge „nichts zu schaffen haben mit diesem Gerechten“, Matth. 27, 19. Pilatus sucht Auswege; er verhängt eine gelindere Strafe, um die Wuth der Juden einigermaßen zu beschwichtigen: er läßt Christum geißeln und wendet sich dann an das Mitleid der Menge. Umsonst! So verurtheilt er denn, indem er seine Hände in Unschuld zu waschen vorgibt, einen unbescholtenen Mann. Das war ein Hohn auf die Gerechtigkeitspflege! Draußen auf dem Richtplatz beichtet der Schächer: „Wir empfangen, was unsere Thaten werth sind; dieser aber hat nichts Ungeschicktes gehandelt“, Luc. 23, 41. Der Hauptmann, der ihn hat verscheiden sehen, spricht ernst und bezeugend: „Fürwahr, dieser ist ein frommer Mensch gewesen“, Luc. 23, 47. Und noch ein stummes Zeugniß der Unschuld Christi: „alles Volk, das dabei war, ... schlugen sich an ihre Brust und wandten wieder um“, Luc. 23, 48.

Ja, die Juden haben, als sie baten, daß man ihnen den Mörder Barabbas schenke, „den Heiligen und Gerechten verleugnet“, Apost. 3, 14. Mit Recht klagt die erste Christengemeinde im Gebet die Mörder Christi an, sie seien versammelt gewesen „über Gottes heiliges Kind Jesum“, Apost. 4, 27. Mit Recht fragt die Kirche noch heute betroffen, erschüttert: „Herzliebster Jesu, was hast du verbrochen?“

Aber eins fällt in der Leidensgeschichte Christi auf: sein Benehmen gegen Gott, seinen Vater. Vor Menschen bezeugt er seine Unschuld und straft sie zuletzt mit dem Schweigen der Unschuld, aber vor Gott windet er sich im Staube wie ein Wurm. Eine unbeschreibliche Angst hat ihn im Garten Gethsemane ergriffen. Er fühlt, daß er Gott wider sich habe, und schreit: „Mein Gott, mein Gott, warum hast du mich verlassen?“ Gott, dem kein gottlos Wesen gefällt und der die Sünder nicht hört, hat sich von Christo abgewendet, als von einem großen Sünder. In Gottes Augen ist der Unschuldige schuldig.

Das ist so gekommen: „Gott hat den, der von keiner Sünde wußte, für uns zur Sünde gemacht“, 2 Cor. 5, 21. Nicht bloß zu einem Sünder, sondern zur Sünde; alle Sünden auf einen Haufen sind dem unschuldigen Christus zugerechnet worden, so daß er in den Augen Gottes dastand als „Sünde“. Die Weltschuld lag auf ihm; er „trug“ sie, Joh. 1, 29. „Der Herr warf unser aller Sünde auf ihn“, Jes. 53, 6.; „er trägt ihre Sünde“, V. 11.; „er hat vieler Sünde getragen“, V. 12. Er hat sie aber auch wahrhaftig als seine Sünden angesehen. „Gott, du weißest meine Thorheit, und meine Schulden sind dir nicht verborgen“, Ps. 69, 6., hat er schon durch Davids Mund geklagt. Und er hat darum auch der Sünden Sold empfangen „an seinem Leibe auf dem Holz“, wohin er unsere Sünden getragen hat, 1 Petr. 2, 24.; er ist „ein Fluch geworden“, Gal. 3, 13. Er war voller Schmerzen und Krankheit, geplagt, geschlagen, gemartert, verwundet, Jes. 53, 3. ff. Die Schuld, wiewohl sie nicht seine eigene, wiewohl sie fremde und Christo zugerechnet war, wurde doch von ihm als seine eigene gefühlt und behandelt. Er bezahlte, was er nicht geraubt hatte, Ps. 69, 5., aber er bezahlte es. Er trank den Kelch, dessen Bitterkeit ihn erbeben machte.

Es mußte so kommen. „Der Herr wollte ihn also zerschlagen“; „des Herrn Vornehmen“ sollte durch Christi Hand „fortgehen“, Jes. 53, 10. Diese Uebertragung fremder Schuld auf einen Unschuldigen, der im Stande und willig war, sie zu tragen, ist „aus bedachtem Rath und Vorsehung Gottes“ geschehen, Apost. 2, 23. Die Mörder Christi, wiewohl sie des Teufels Gesellen waren, haben gethan, „was Gottes Hand und Gottes Rath zuvor bedacht hat, daß geschehen sollte“, Apost. 4, 28. Die Schlange biß nach der Ferse des Weibesamens, und Gott ließ das geschehen, damit eben dann ihr Haupt zermalmt würde. Er hat die Hölle in ihr Verderben geschickt; denn an dem, den sie verschlingen wollte, hat sie sich den Tod und die Pestilenz geholt, Hos. 13, 14.

Aber nicht an Christo allein hat Satan verspielt, sondern auch an uns. \gesperrt{Christi Sieg ist uns gegeben}. Die Uebertragung unserer Schuld auf den Unschuldigen ist nur eine Seite des wunderlichen Wechsels, den St. Paulus 2 Cor. 5, 21. nennt. Die andere Seite heißt: „auf daß \gesperrt{wir} würden in ihm \gesperrt{die Gerechtigkeit, die vor Gott gilt}“. Gott nahm von uns, was wir hatten, und gab Christo, was er nicht hatte: Sünde; Gott nahm von Christo, was er hatte, und gab uns, was wir nicht hatten: Gerechtigkeit. Wir haben gesehen, daß die erste Seite dieses Tausches kein Spiel, keine Einbildung, keine bloße Uebertragung in Gedanken, sondern wirklicher, wahrhaftiger, ja, bitterer Ernst war. Der zweite Theil des wunderlichen Wechsels ist ebenso wirklich und wahrhaftig. Das Blut Jesu Christi, des Sohnes Gottes, „macht uns rein von aller Sünde“, 1 Joh. 1, 7. „Die Strafe liegt auf ihm, auf daß wir Friede hätten, und durch seine Wunden sind wir geheilet“, Jes. 53, 5. Dieser „Knecht Gottes, der Gerechte, macht viele gerecht; denn er trägt ihre Sünde“, Jes. 53, 11. „Gott war in Christo und versöhnte die Welt mit ihm selber, und rechnete ihnen ihre Sünden nicht zu“, 2 Cor. 5, 19. Wir mit unserem eitlen Wandel nach väterlicher Weise, wir Unreinen, wir Kinder vieler Missethaten sind durch die Barmherzigkeit Gottes und das Blut seines Sohnes Heilige und Gerechte. Wir sind nicht einigermaßen gerecht, halbwegs schuldlos, sondern wir sind in Christo „die Gerechtigkeit, die vor Gott gilt“. Wir sind nun ohne alle und jede Schuld; wir können in der Gerechtigkeit Christi vor Gott bestehen. An denen, die in Christo sind, ist „nichts Verdammliches“, Röm. 8, 1. Ihr Gewissen beißt sie nicht mehr; ihr Herz verdammt sie nicht, 1 Joh. 3, 21.; sie kommen nicht in das Gericht, Joh. 5, 24. Sie haben durch den Glauben den Tausch, welchen Gottes Gnade ihnen zu ihrem Heile im Wort der Versöhnung anbietet, angenommen, und Gott betrachtet, erklärt und behandelt sie nun als Gerechte. Die Schuldigen sind schuldlos geworden.

Fürwahr, „er wechselt mit uns wunderlich“.

W. H. T. D.
